\chapter{DATA ANALYSIS -- INDUSTRIAL RECONFIGURATION}
\label{ch:industry_reconfig}

\section{Introduction}
\label{sec:reconfig_intro}

This chapter presents the analysis of the supply-side data
collected from 50 industry specialists (33 structured-only, 17
structured + interview, of whom 28 were transcribed in full).  It
addresses Research Question~3, tests Hypotheses H6--H8 and provides
the qualitative interpretive frame for the demand-side findings of
Chapter~\ref{ch:viewership}.  The chapter is organised by the eight
substantive specialist-side figures (Figures~\ref{fig:specialist_roles}
through~\ref{fig:marketing}), supplemented by a thematic-coding
synthesis of the interview data (Section~\ref{sec:themes}) and a
hypothesis-test summary (Section~\ref{sec:hyp_summary_6}).

Quantitative data from the specialist instrument were tabulated in
SPSS 28 with one-sample $t$-tests against pre-pandemic baselines
where applicable.  Qualitative interview data were coded in NVivo
14 using the constant-comparative approach described in
Section~\ref{sec:analysis}.

\section{Profile of the Specialist Sample}
\label{sec:specialist_profile}

The role distribution of the 50-respondent specialist sample is
shown in Figure~\ref{fig:specialist_roles}.  The sample spans the
Tamil-cinema value chain: producers (10), directors (8), OTT
acquisition / commissioning heads (7), distributors (7), exhibitors
(6), marketing / PR professionals (5), music-label executives (4)
and trade analysts (3).

\begin{figure}[h]
\centering
\includegraphics[width=0.88\textwidth]{fig_21_specialist_roles.png}
\caption{Composition of industry-specialist sample.}
\label{fig:specialist_roles}
\end{figure}

Years of experience in the industry are presented in
Figure~\ref{fig:specialist_exp}.  Sixty percent of respondents have
between 10 and 20 years' experience and a further 18\% have over 20
years', supporting the interpretive weight of the specialist
narratives.

\begin{figure}[h]
\centering
\includegraphics[width=0.62\textwidth]{fig_22_specialist_exp.png}
\caption{Years of industry experience.}
\label{fig:specialist_exp}
\end{figure}

Of the 50 specialists, 17 are women (34\%) -- a higher share than the
all-India industry average reported by \citet{ficciey2024} but
consistent with the rising representation of women in OTT
commissioning, marketing and trade-press functions
\citep{kpmg2023,mukherjee2020}.

\section{Specialist-Perceived COVID Impact on Revenue Streams}
\label{sec:revenue_impact}

Specialists were asked to estimate, on the basis of their own
projects, the median percentage change in each of six revenue
streams between the pre-pandemic baseline (CY 2019) and the
post-pandemic equilibrium (CY 2024).  The median responses are
plotted in Figure~\ref{fig:revenue_impact}.

\begin{figure}[h]
\centering
\includegraphics[width=0.92\textwidth]{fig_23_revenue_impact.png}
\caption{Specialist-perceived COVID impact by revenue stream.}
\label{fig:revenue_impact}
\end{figure}

The clearest patterns are: (i) a 58\% median decline in theatrical
revenue per film, partly reflecting capacity contraction (Table~\ref{tab:exhibition},
Chapter~\ref{ch:industry_overview}) and partly the redirection of
viewer time to OTT; (ii) a 47\% increase in OTT licensing revenue,
the only stream to grow at scale; (iii) a 19\% decline in satellite
TV revenue, reflecting the secular substitution of satellite by
streaming; (iv) a 12\% rise in music / digital rights revenue,
driven by streaming-audio royalties; (v) a 33\% decline in
international (overseas) box-office revenue, driven by more permissive
diaspora OTT options; and (vi) a 64\% decline in in-cinema food and
beverage (F\&B) revenue per release, magnifying the theatrical
revenue contraction at the exhibitor level.

\textbf{Implication.}  No single revenue stream is healthier than its
2019 baseline except OTT.  This finding is consistent with the
demand-side spend-mix shift documented in Section~\ref{sec:spend_results}
of Chapter~\ref{ch:viewership}.

\section{Theatrical-Window Negotiation}
\label{sec:windowing_findings}

Specialists reported the distribution of theatrical-to-OTT release
windows for the projects they had been involved in across two
periods: the pre-COVID baseline (films released in 2018--2019) and
the post-COVID period (films released in 2023--2024).  The
distributions are shown in Figure~\ref{fig:window_negotiated}.

\begin{figure}[h]
\centering
\includegraphics[width=0.85\textwidth]{fig_24_window_negotiated.png}
\caption{Theatrical-to-OTT window negotiated by specialists.}
\label{fig:window_negotiated}
\end{figure}

The pre-COVID modal negotiated window was 5--8 weeks (56\% of
films), with a residual long-tail above 8 weeks (26\%).  Post-COVID,
the modal window has shifted to 3--4 weeks (41\%), with a meaningful
short-window block at 1--2 weeks (22\%) and a small but
non-negligible 6\% of releases negotiated for same-day OTT
availability.

A chi-square test of homogeneity on the underlying frequency
distribution is highly significant: $\chi^2(4) = 53.2$, $p < 0.001$.
H6 is supported.  The supply-side window has shortened decisively
and the shortening is bi-modal: most films cluster in the 3--4 week
band, but a meaningful sub-set has migrated to near- or simultaneous
release.

The viewer-side preference (Section~\ref{sec:window_results}) and the
supply-side negotiated reality have therefore converged on the
3--4 week window as the new modal arrangement.  This is one of the
strongest findings of the thesis and is given its own discussion in
Section~\ref{sec:hybrid_eq} of Chapter~\ref{ch:conclusion}.

\section{OTT Acquisition Price Index}
\label{sec:price_index}

Figure~\ref{fig:ott_price} traces the specialist-reported OTT
acquisition price index across three film tiers, indexed to the
2018 baseline ($=100$).

\begin{figure}[h]
\centering
\includegraphics[width=0.88\textwidth]{fig_25_ott_price.png}
\caption{Specialist-reported OTT acquisition-price index.}
\label{fig:ott_price}
\end{figure}

The trajectory shows a sharp 2020--2021 spike at every tier, peaking
at 244 (star-led), 192 (mid-budget) and 138 (small / indie) before
correcting in 2023--2024 as platform commissioning budgets were
re-disciplined and competition for distinctive titles intensified
\citep{ormax2024,kpmg2023,mpa2024}.  The 2025 index value remains
above 2018 baseline at the star (156) and mid-budget (108) tiers but
\emph{below} baseline at the small / indie tier (78).

A one-sample $t$-test of the 2025 specialist-reported acquisition
price (per film) against the 2018-equivalent (deflated to constant
INR using CPI; \citealp{rbi2024}) yields, for star-led films,
$t = 4.7$, $p < 0.001$, indicating the 2025 price remains
\emph{statistically significantly above} the 2018 baseline.

\textbf{Implication.}  H8 is supported.  The acquisition-price market
has corrected from the 2020--2021 over-shoot but has not returned to
the pre-pandemic level for star-led titles -- evidence that OTT
licensing has become a structurally larger income stream rather than
a temporary revenue substitute.

\section{Production-Budget Restructuring}
\label{sec:budget_results}

Figure~\ref{fig:budgets} compares the distribution of Tamil films
greenlit in the pre-COVID baseline (2018--2019) with the post-COVID
period (2023--2024) by production-budget bracket.

\begin{figure}[h]
\centering
\includegraphics[width=0.85\textwidth]{fig_26_budgets.png}
\caption{Production-budget bracket distribution.}
\label{fig:budgets}
\end{figure}

The distribution exhibits compression toward the lower budget
brackets.  The proportion of films in the under-INR-5-crore bracket
rose from 22\% to 34\%; the 5--15 crore bracket from 38\% to 41\%;
and the upper brackets (above 15 crore) contracted, with the
30--60 crore band roughly halving from 11\% to 7\%.  The over-60
crore band has remained small but is now constituted by a
recognisably different mix of titles -- pan-India star-led
spectacles such as the Sivakarthikeyan / Kamal Haasan / Vijay /
Suriya tier -- with multiplex- and OTT-targeted economics.

\textbf{Implication.}  The mid- and large-budget contraction is
consistent with the specialist-narrated risk-aversion documented in
the qualitative interview data (Section~\ref{sec:themes}) and with
the \citet{caves2000} ``nobody knows'' principle of cultural-industry
risk management.  Producers are choosing to greenlight a higher
volume of smaller-budget, content-driven films.

\section{Revenue-Mix Reconfiguration of an Average Tamil Release}
\label{sec:revenue_mix_results}

Figure~\ref{fig:revenue_mix} presents the specialist-reported
revenue-mix of an ``average Tamil release'' in the pre- and post-COVID
periods.

\begin{figure}[h]
\centering
\includegraphics[width=0.85\textwidth]{fig_27_revenue_mix.png}
\caption{Revenue-mix reconfiguration of an average Tamil release.}
\label{fig:revenue_mix}
\end{figure}

The theatrical share fell from 62\% to 44\%; the OTT share rose
from 6\% to 28\%; satellite fell from 17\% to 13\%; music / digital
rose marginally from 8\% to 9\%; and international fell from 7\% to
6\%.  A one-sample $t$-test of the post-COVID OTT share
($M = 28\%$, $SD = 7$) against the pre-COVID baseline (6\%) is
highly significant: $t(49) = 22.2$, $p < 0.001$.  H7 is supported.

\textbf{Implication.}  OTT has migrated from a peripheral, post-theatrical
windowed revenue stream to a structurally central component of the
Tamil-film commercial proposition.  The combined ``digital'' share
(OTT + music / digital) rose from 14\% to 37\%, indicating that
``digital'' broadly construed has become the second leg of a
two-legged revenue model alongside theatrical.

\section{Star-Fee Restructuring}
\label{sec:star_fees_results}

Figure~\ref{fig:star_fees} reports the specialist-estimated median
fee adjustment between 2019 and 2024 by star tier.

\begin{figure}[h]
\centering
\includegraphics[width=0.85\textwidth]{fig_28_star_fees.png}
\caption{Specialist-reported star-fee restructuring (2019$\rightarrow$2024).}
\label{fig:star_fees}
\end{figure}

The findings are theoretically pivotal.  The top tier
(``Vijay / Ajith / Kamal / Rajini-class'') has experienced a
relatively modest median fee correction of $-12\%$.  The A-list
(immediately below the top tier) has corrected by $-22\%$, the B-list
by $-34\%$, and character / supporting actors by $-18\%$.  Because of
small-sample variability across specialists, we report 95\% confidence
intervals: top-tier $[-18\%,-6\%]$; A-list $[-29\%,-15\%]$; B-list
$[-42\%,-26\%]$; character $[-26\%,-10\%]$.

The pattern admits a coherent reading.  The pure ``top-tier'' star
remains a near-uniqueness in the Tamil market, with continued
opening-day mobilisation power that protects fee.  The A- and
B-list -- which previously could command large fees because of
default star-driven commissioning logic -- have corrected most
sharply because the demand side has migrated toward content-driven
choice (Section~\ref{sec:star_content_results} of
Chapter~\ref{ch:viewership}).  Character / supporting fees have
fallen modestly because budget compression
(Section~\ref{sec:budget_results}) has affected all line items.

\section{Marketing-Channel Adoption}
\label{sec:marketing_results}

Figure~\ref{fig:marketing} reports the percentage of specialists
who report using each marketing channel for ``most or all'' of
their releases.

\begin{figure}[h]
\centering
\includegraphics[width=0.92\textwidth]{fig_29_marketing.png}
\caption{Adoption of marketing channels (\% of specialists).}
\label{fig:marketing}
\end{figure}

The most striking shifts are: (i) influencer marketing rising from
28\% to 78\%; (ii) OTT in-app promotion from 14\% to 72\%; (iii)
YouTube digital premiere from 24\% to 68\%; and (iv) the modest
\emph{decline} of theatrical hoardings from 88\% to 64\%.  The
audio-launch live-stream rose from 36\% to 62\%, indicating that
the hybridisation of the audio-launch event has been industrialised
rather than displaced.

\textbf{Implication.}  Tamil-cinema marketing has migrated from a
predominantly theatre-and-music-centric stack to a digital-first
stack.  The theatrical hoarding -- the iconic visual of mid-Chennai
-- has not disappeared but has lost its share of mind to digital
channels.

\section{Qualitative Themes from Specialist Interviews}
\label{sec:themes}

The 28 transcribed interviews were coded in NVivo 14 through three
cycles of coding as described in Section~\ref{sec:analysis}.  The
nine themes identified are reported in Table~\ref{tab:themes} along
with their prevalence and a representative anonymised quotation
each.  The names of specialists are withheld; designations are
generic.

\begin{table}[h]
\centering
\caption{Thematic coding of specialist interviews (n=28)}
\label{tab:themes}
\small
\begin{tabular}{p{3.5cm} p{1.4cm} p{9.5cm}}
\toprule
\textbf{Theme} & \textbf{Cited by} & \textbf{Representative quotation (paraphrased)}\\
\midrule
Risk re-assessment of star-only logic & 24/28 &
``We assumed for two decades that a star meant safety.  After 2020
we learned a star can also mean a 90-crore liability if the script
isn't right.''  --- Senior producer.\\
Hybrid window as new normal & 26/28 &
``Four weeks is the working compromise.  Less than four, the
exhibitors revolt; more than four, the OTT pulls back the price.''
--- OTT acquisition head.\\
Mid-budget renaissance & 22/28 &
``Eight to twelve crore, with a clean script and a known director,
is the safest commercial proposition in the industry today.''
--- Producer.\\
Direct-to-OTT as portfolio play & 19/28 &
``We don't see direct-to-OTT as failure any more; we structure
it as a different kind of release.''  --- Marketing head.\\
Genre experimentation & 23/28 &
``OTT taught the audience to like content they would have walked
out of in a theatre five years ago.''  --- Director.\\
Marketing convergence on digital & 27/28 &
``Influencer placements convert better and cost less than full-page
ads.''  --- PR specialist.\\
Workforce informalisation & 17/28 &
``A lot of mid-tier technicians simply did not return after the
lockdown; the industry is younger and more freelance now.''
--- Studio head.\\
Mid-tier star anxiety & 21/28 &
``The A-list has felt the price-correction the most; the top tier
is insulated, the bottom tier was always priced low.''
--- Trade analyst.\\
Pan-India ambition & 20/28 &
``Every meaningful pitch in 2024 began with -- can this be a pan-India
release?''  --- Producer.\\
\bottomrule
\end{tabular}
\end{table}

The nine themes map directly onto the four-quadrant Hybrid-Equilibrium
framework presented in Chapter~\ref{ch:conclusion}: themes 1, 2 and
3 anchor the ``hybrid distribution'' quadrant; themes 5 and 8 anchor
the ``content-led demand'' quadrant; themes 3 and 9 anchor the
``mid-budget central unit'' quadrant; and themes 4 and 6 anchor the
``OTT-native production'' quadrant.

\section{Hypothesis Test Summary (H6--H8)}
\label{sec:hyp_summary_6}

\begin{table}[h]
\centering
\caption{Hypothesis test summary -- specialist-side hypotheses}
\label{tab:hyp_6}
\small
\begin{tabular}{p{4.0cm} p{4.0cm} p{2.4cm} p{2.0cm} p{1.5cm}}
\toprule
\textbf{Hypothesis} & \textbf{Test} & \textbf{Stat} & \textbf{$p$} & \textbf{Status}\\
\midrule
H6: 5--8 wk window share fell & Chi-square homogeneity &
$\chi^2 = 53.2$ & $<$0.001 & Supported \\
H7: OTT revenue share rose    & One-sample $t$ vs 6\% baseline & $t = 22.2$ & $<$0.001 & Supported \\
H8: 2024 acquisition price corrected but $>$ 2018 (star tier) & One-sample $t$ vs 2018 deflated baseline & $t = 4.7$ & $<$0.001 & Supported \\
\bottomrule
\end{tabular}
\end{table}

All three supply-side hypotheses are supported.  Combined with the
five viewer-side hypotheses of Chapter~\ref{ch:viewership}, this
gives seven of eight hypotheses supported and one (H5) partially
supported.

\section{Synthesis: A Recombinatory Industry}
\label{sec:synthesis_6}

The supply-side findings, taken together, do not describe an industry
that has been disrupted in the Christensenian sense
\citep{christensen1997}.  They describe an industry that has
\emph{recombined} -- one that has absorbed the streaming entrant into
its existing value chain, restructured its windowing logic,
recalibrated its star economics, compressed its budget brackets and
re-engineered its marketing stack.  The result is a leaner,
mid-budget, content-led, hybrid-distribution Tamil-cinema industry
that is theoretically interesting precisely because it has avoided
the binary substitution that early commentators predicted
\citep{wayne2018,fagerjord2018}.

The qualitative interview data corroborate the quantitative
distributions at every point of contact: the windowing data
(Section~\ref{sec:windowing_findings}) is supported by Theme 2; the
star-fee data (Section~\ref{sec:star_fees_results}) by Theme 8;
the budget data (Section~\ref{sec:budget_results}) by Theme 3; the
marketing data (Section~\ref{sec:marketing_results}) by Theme 6.
The convergence between the two methodological strands is itself
the strongest evidence for the validity of the joint argument.

\section{Chapter Summary}
\label{sec:reconfig_summary}

This chapter has presented the analysis of the 50-respondent
specialist sample, addressing Research Question~3 and testing
hypotheses H6--H8.  The supply-side findings document a Tamil film
industry that has materially reconfigured its windowing, budget,
star-fee, marketing and revenue-mix arrangements in response to the
demand-side shifts documented in Chapter~\ref{ch:viewership}.  Nine
qualitative themes corroborate the quantitative findings.  All
three supply-side hypotheses are supported.  The next chapter
synthesises the demand- and supply-side findings into the Hybrid-Equilibrium
framework and derives theoretical, managerial and policy implications.
